\documentclass[11pt]{article}
\usepackage{amsmath, amssymb, amsthm}
\usepackage[a4paper,margin=2.5cm]{geometry}
\newcommand{\Poi}{\mathrm{Poi}}
\newcommand{\NN}{\mathbb{N}_0}
\newcommand{\E}{\mathbb{E}}
\begin{document}

\section*{Two Wright--Fisher updates and their stationary marginals}

\subsection*{Setup}
Each individual carries a mutation count $k\in\NN$. We use exponential
selection $w_k = e^{-\alpha k}$ ($\alpha>0$) and Poisson mutation at rate
$\lambda\ge 0$ (per individual, per generation), so the new mutation count of an
offspring is the parent's class plus an independent $\Poi(\lambda)$.

Let $Q$ denote a probability distribution on $\NN$ (the marginal class
distribution of one individual). Two basic operators act on such $Q$:
\begin{align*}
  (M_\lambda Q)(k) &\;=\;(Q\ast\Poi(\lambda))(k)
                \;=\;\sum_{j=0}^{k} Q(j)\,\Poi(\lambda;k-j)
                \qquad\text{(mutation)},\\
  (S\,Q)(k)     &\;=\;\frac{w_k\,Q(k)}{\sum_{j\ge 0} w_j\,Q(j)}
                \qquad\text{(selection / reweighting by $w$).}
\end{align*}
We compare two update rules per generation:
\[
  \text{(A) ``sel $\to$ mut''}:\quad T_A := M_\mu\circ S,
  \qquad
  \text{(B) ``mut $\to$ sel''}:\quad T_B := S\circ M_\lambda.
\]
Multinomial resampling of $N$ offspring is applied afterwards; for the
analysis of the marginal we look at the deterministic large-$N$ limit and
return to fluctuations at the end.

\subsection*{Two Poisson identities}
The reason the calculation is short is that both operators preserve the
Poisson family:
\begin{align}
  S\bigl(\Poi(\nu)\bigr)         &\;=\;\Poi\bigl(\nu\,e^{-\alpha}\bigr),
  \label{eq:sel}\\
  M_\lambda\bigl(\Poi(\nu)\bigr)    &\;=\;\Poi\bigl(\nu+\lambda\bigr).
  \label{eq:mut}
\end{align}
(Eq.~\eqref{eq:sel} follows from
$w_k\,\Poi(\nu;k)=e^{-\alpha k}e^{-\nu}\nu^k/k!\propto(\nu e^{-\alpha})^k/k!$,
i.e.~exponentially tilting Poisson rescales its parameter.
Eq.~\eqref{eq:mut} is the standard sum-of-independent-Poissons identity.)

\subsection*{Stationary Poisson marginals}
Suppose the marginal is $Q=\Poi(\theta_{\mathrm{eq}})$ and we look for
fixed points of $T_A,T_B$.

\paragraph{Model A (sel $\to$ mut).}
By \eqref{eq:sel} and \eqref{eq:mut},
\[
  T_A\bigl(\Poi(\theta_{\mathrm{eq}})\bigr)
  = M_\mu\bigl(\Poi(\theta_{\mathrm{eq}}e^{-\alpha})\bigr)
  = \Poi\bigl(\theta_{\mathrm{eq}}\,e^{-\alpha}+\mu\bigr).
\]
The fixed-point equation
$\theta_{\mathrm{eq}}=\theta_{\mathrm{eq}}\,e^{-\alpha}+\mu$ gives
\begin{equation}\label{eq:thetaA}
  \boxed{\;\theta_{\mathrm{eq}}^{(A)}\;=\;\frac{\mu}{1-e^{-\alpha}}\;.}
\end{equation}

\paragraph{Model B (mut $\to$ sel).}
Now
\[
  T_B\bigl(\Poi(\theta_{\mathrm{eq}})\bigr)
  = S\bigl(\Poi(\theta_{\mathrm{eq}}+\lambda)\bigr)
  = \Poi\bigl((\theta_{\mathrm{eq}}+\lambda)\,e^{-\alpha}\bigr),
\]
so $\theta_{\mathrm{eq}}=(\theta_{\mathrm{eq}}+\lambda)e^{-\alpha}$ yields
\begin{equation}\label{eq:thetaB}
  \boxed{\;\theta_{\mathrm{eq}}^{(B)}
        \;=\;\frac{\lambda\,e^{-\alpha}}{1-e^{-\alpha}}
        \;=\;\frac{\lambda}{e^{\alpha}-1}\;.}
\end{equation}

For small $\alpha$, both formulas reduce to $\theta_{\mathrm{eq}}\approx\lambda/\alpha$, but for
finite $\alpha$ they differ; in particular
\[
  \frac{\theta_{\mathrm{eq}}^{(B)}}{\theta_{\mathrm{eq}}^{(A)}}
  \;=\;\frac{\lambda}{\mu}\cdot\frac{1-e^{-\alpha}}{e^{\alpha}-1}
  \;=\;\frac{\lambda}{\mu}\cdot e^{-\alpha}.
\]

\end{document}
